\documentclass[11pt]{article}
\usepackage[T1]{fontenc}
\usepackage[utf8]{inputenc}
\usepackage[english]{babel}
\usepackage[margin=1in]{geometry}
\usepackage{amsmath,amssymb}
\usepackage{enumitem}
\usepackage{microtype}

\ifdefined\pdfcatalog
  \pdfcatalog{/Lang (en-US)}
\fi

\setlength{\parindent}{0pt}
\setlength{\parskip}{0.55em}
\setlength{\emergencystretch}{2em}
\setlist{nosep,leftmargin=*}

\newcommand{\findinglabel}[1]{\textbf{#1}}
\newcommand{\classification}[1]{\textit{Classification: #1.}}

\title{Technical Review of ``Early-Stopping Activation Steering for Factual Preference Alignment in Vietnamese Medical Language Models''}
\author{Manuscript review}
\date{August 28, 2026}

\begin{document}
\selectlanguage{english}
\maketitle

\section*{Overall assessment}

This revision resolves several definite faults in earlier drafts.  The hook target is consistently Layer~8, the 100-vector placebo summary is arithmetically coherent, the natural Hard-Cutoff confidence interval is centered on its 5.20-point effect, the source completes a two-pass build, and the RAG comparison includes an oracle condition, retrieval parameters, paired discordant counts, and latency dispersion.  In particular, the reported RAG-versus-steering McNemar calculation is consistent with $b=21$, $c=64$, and an 8.60-point paired difference.

The manuscript still requires major revision.  Its bounded-generation paragraph points to a natural-completion table that does not contain the claimed $2\times2$ outcomes, while the natural tests remain unreproducible because their discordant counts and test variants are absent.  The placebo result is presented simultaneously as a $5.15\sigma$ effect and a finite-resolution empirical $p=0.0099$, which are not equivalent inferential statements.  The activation section is more cautious, but the Abstract and contribution list still claim saturation prevention and stabilization from scalar norms.  RefPref remains a semantic proxy, the natural experiment remains a trade-off rather than a uniform quality gain, and dataset, metric, timing, vector, hook, and clinical-validation details remain incomplete.  The natural-completion table also overflows its column by approximately 61.44~pt.

\section*{Major technical findings}

\subsection*{1. The bounded and natural paired analyses remain structurally mixed and incompletely reproducible}

\classification{Definite cross-reference error and statistical ambiguity}

\findinglabel{Location.} Controlled Bounded-Generation Stress Test; paragraph preceding Table~\texttt{tab:mcnemar}; Table~\texttt{tab:mcnemar}; Natural Termination Validation.

\findinglabel{Problem.} The bounded-generation paragraph states that Table~\texttt{tab:mcnemar} reports the paired $2\times2$ outcomes and then gives $b=16$, $c=37$, net change $21/500=4.20$~pp, and exact $p=0.0055$.  The referenced table instead contains three natural-completion schedule comparisons at \texttt{max\_new\_tokens=800}; it contains neither the bounded $2\times2$ table nor any discordant counts.

The revised natural Hard-Cutoff row is numerically plausible: a 26-record net change and $p\approx0.00086$ are compatible, for example, with $b=16$, $c=42$.  However, the manuscript never reports those counts.  It also does not identify whether the Continuous and Linear-Decay values are exact or asymptotic McNemar tests.  The reported $p=0.0412$ for Continuous is compatible with a continuity-corrected asymptotic calculation for one possible transition table, illustrating why the missing test definition matters.  Bootstrap unit, repetitions, seed, and interval construction are also absent.

\findinglabel{Why it matters.} A reader cannot reconstruct the primary bounded table or verify the natural paired tests and multiplicity adjustment from the marginal accuracies alone.

\findinglabel{Suggested fix.} Provide separate bounded and natural tables.  For every comparison, report all four paired cells, identify $b$ and $c$, state exact versus asymptotic testing, give the software call, and document the bootstrap procedure and Holm ordering.  Generate the prose, captions, CIs, and adjusted values from the same versioned analysis artifact.

\subsection*{2. The placebo evidence is useful but the $\sigma$ interpretation is overstated}

\classification{Unsupported strength of statistical interpretation}

\findinglabel{Location.} Abstract; contribution~4; placebo Methods and Results; Table~\texttt{tab:baselines}; Conclusion.

\findinglabel{Problem.} The updated quantities are internally consistent:
\[
\frac{77.20-55.82}{4.15}\approx5.15,
\qquad
\frac{1}{100+1}\approx0.0099.
\]
The first value is a standardized distance using the observed dispersion of 100 bounded accuracy values.  It is not automatically a calibrated Gaussian ``five-sigma'' significance statement, especially because accuracy is bounded and normality is not assessed.  The second is the smallest plus-one empirical randomization value obtainable from 100 directions when none exceeds the target.  Thus $5.15\sigma$ and empirical $p=0.0099$ have very different evidential meanings.  The paper does not define whether $4.15\%$ is a sample SD, SE, or CI, and the table supplies one placebo BERTScore value, 0.6945, without dispersion across directions.

\findinglabel{Why it matters.} The controls support direction specificity, but the present language can be read as conventional five-sigma evidence even though the randomization test has only 0.0099 resolution.

\findinglabel{Suggested fix.} Call 5.15 a descriptive standardized distance, define the $\pm$ quantity, and report the empirical result as ``0 of 100 sampled directions exceeded the target; plus-one $p=1/101$.''  Release all per-direction RefPref and BERTScore values and seeds.  Use more sampled directions if a smaller randomization value is important, and avoid converting the empirical result into Gaussian-tail language.

\subsection*{3. Scalar norm dynamics still do not demonstrate saturation prevention}

\classification{Unsupported mechanism claim and notation drift}

\findinglabel{Location.} Abstract; Introduction; contribution~3; Early-Stopping Steering Hook \& Saturation Diagnostics; Empirical Activation Magnitude Dynamics.

\findinglabel{Problem.} The Results correctly calls 51.20 a $-3.01$-unit undershoot and explicitly cautions that norms do not capture manifold geometry.  The Abstract nevertheless says activation tracking ``confirms'' the mechanism and calls the undershoot stabilization; contribution~3 remains titled ``Empirical Saturation Diagnostics.''  At $t=50$, the Linear-Decay norm is farther from the 54.21 baseline in absolute value than the Continuous norm is:
\[
|51.20-54.21|=3.01,
\qquad
|56.18-54.21|=1.97.
\]
No acceptable baseline band or behavioral stability criterion is defined.

Only two selected time points are given, without dispersion, confidence bands, per-step survivor counts, or association with repetition and factual errors.  For $t>K$, early-stop conditions have no direct perturbation, but their free-generation histories contain different preceding tokens; norm differences therefore conflate intervention effects with changed contexts.  Notation also drifts: the Abstract and contribution use $\|h_8^{(t)}\|_2$ while calling the value post-hook, whereas Results uses $\|\hat h_8^{(t)}\|_2$.

\findinglabel{Why it matters.} The central mechanistic claim remains stronger than the diagnostic, which establishes only post-hook norm elevation and undershoot under different generation trajectories.

\findinglabel{Suggested fix.} Define saturation operationally.  Add fixed-token replay or teacher forcing, full pre/post-hook trajectories, uncertainty and survivor counts, projection onto $v_{\text{steer}}$, cosine drift, covariance or distributional distance, and response to an additional controlled perturbation.  Relate these diagnostics to output errors.  Until then, use ``activation-norm dynamics'' consistently and remove ``confirms saturation prevention'' and ``stabilization.''

\subsection*{4. Natural completion remains a metric trade-off and does not establish Linear Decay as the best schedule}

\classification{Unsupported comparative conclusion}

\findinglabel{Location.} Abstract; contribution~2; natural-completion Results; Table~\texttt{tab:long\_gen}; Discussion; Conclusion.

\findinglabel{Problem.} All steering schedules raise RefPref but lower reference BERTScore and increase Rep-4 relative to baseline.  Hard Cutoff has the highest RefPref (70.60\%), baseline has the highest BERTScore (0.6779) and lowest Rep-4 (4.10\%), and Linear Decay has the highest repetition (5.37\%).  Hard Cutoff's Rep-4 increase to 5.35\% is 30.5\% relative, so the Results description ``minor increase'' is not justified without uncertainty or a practical threshold.  The manuscript's own natural table labels Linear Decay versus baseline nonsignificant at $p=0.1250$, yet the Conclusion states that linear decay yields consistent improvements.

Contribution~2 and the opening natural-completion paragraph still say EOS Hit is 100.00\%, contradicting the condition-specific 99.80\% values for Continuous Steering and Linear Decay.  In addition, the proposed paper emphasizes linear decay, while the strongest natural RefPref result belongs to Hard Cutoff.

\findinglabel{Why it matters.} The data do not demonstrate uniform output-quality improvement or establish the proposed decay schedule as the preferred natural-generation intervention.

\findinglabel{Suggested fix.} Prespecify one primary endpoint and schedule on validation data and evaluate that frozen choice once on test data.  Report paired CIs for RefPref, BERTScore, and Rep-4, with noninferiority margins if claiming acceptable degradation.  Use the 99.80\%--100.00\% EOS range and metric-specific comparative language throughout.

\subsection*{5. RefPref is not a clinical factual-correctness or hallucination endpoint}

\classification{Unsupported interpretation of a disclosed proxy}

\findinglabel{Location.} title; Abstract; Introduction; contributions; Results; Discussion; Conclusion; keywords.

\findinglabel{Problem.} The Discussion identifies RefPref as a semantic proxy, but the explicit Methods limitation in a preceding draft has been removed.  Elsewhere the vector is called ``truthfulness,'' RefPref is repeatedly called ``accuracy'' or ``factual alignment,'' and the paper is framed as hallucination mitigation and clinical decision support.  Relative BERTScore similarity to $y^+$ rather than one synthetic $y^-$ cannot reliably detect a wrong dose, unit, negation, contraindication, or interaction mechanism.  The natural results themselves show RefPref increasing while similarity to the reference decreases.

Equation~4 also maps ties to zero through the strict indicator, while the prose calls ties ``neither reference-preferred nor hallucination-preferred.''  The paper does not state whether ties remain in the denominator as RefPref failures, are excluded, or form a third category in paired tests.

\findinglabel{Why it matters.} Clinical safety, harmful-error reduction, and deployment suitability cannot be inferred from the current metric, and ambiguous tie handling can change counts and paired inference.

\findinglabel{Suggested fix.} Restore the limitation next to the metric definition and specify tie counts, denominator, and paired-test handling.  Add blinded clinician-adjudicated correctness and harmful-error rates, stratified by the three clinical categories, with reviewer qualifications, agreement, and adjudication.  Otherwise narrow the title, Abstract, keywords, and conclusions to reference preference.

\subsection*{6. The RAG study is stronger but still does not support a broad RAG comparison or all efficiency claims}

\classification{Likely weak baseline and incomplete systems interpretation}

\findinglabel{Location.} Abstract; Experimental Setup; RAG Results; Table~\texttt{tab:baselines}; Conclusion.

\findinglabel{Problem.} The revision appropriately adds BM25 parameters, word segmentation, oracle results, and a reproducible paired comparison: $b=21$, $c=64$ implies $43/500=8.60$~pp and exact $p\approx3.28\times10^{-6}$.  However, BM25 Recall@1 remains only 19.20\%.  Passage construction, query formulation, relevance mapping, and retrieval-failure handling are not specified, and no stronger lexical or dense retriever is evaluated.  Oracle RAG reaches 89.40\%, substantially above steering's 77.20\%, indicating that the observed non-oracle deficit is primarily a retrieval-system result rather than evidence against RAG as a method class.

The $\pm$ latency terms are not defined as across-query SDs, repeated-run SDs, SEs, or CIs.  Retrieval, tokenization, prefill, and decoding time are not separated.  Prompt-token counts imply context overhead, but the ``zero-context-memory advantage'' claim is not backed by measured peak memory or KV-cache allocation.  The steering-versus-baseline timing row repeats exactly $7.32\pm0.84$ for multiple conditions, without an equivalence test or independent measurement description.

\findinglabel{Why it matters.} The paper supports superiority to one specified low-recall BM25 configuration, not to RAG generally, and latency or memory advantages depend on measurement boundaries.

\findinglabel{Suggested fix.} Publish passage construction, queries, relevance labels, and failure handling; add stronger retrievers and retrieval-conditioned generation analysis.  State the scope as one non-oracle BM25 setup.  Define the latency statistic, run repeated component-wise timing, measure peak memory, and use an equivalence margin for the claim of no steering overhead.

\subsection*{7. Dataset construction and experimental selection remain insufficiently documented}

\classification{Reproducibility limitation and likely leakage risk}

\findinglabel{Location.} ViHaluEval-Medical Benchmark Construction; Experimental Setup \& Reproducibility Environment; layer and hyperparameter claims.

\findinglabel{Problem.} ``Expert-structured'' remains undefined.  The manuscript provides no annotation protocol, counter-answer generation rules, clinical review fraction, agreement, adjudication, deduplication, source identifiers, or licensing.  The category counts 165, 160, and 175 sum to the 500-question evaluation subset but are placed in the paragraph describing all 14,700 records, so the full-dataset category distribution is unclear.  Question-disjoint splitting does not rule out drug/source-entry or error-template leakage, and the 500-item selection rule and seed are absent.

The environment paragraph supplies useful software versions and placebo seeds, but it does not state how Layer~8, $\alpha_0=18$, and $K=16$ were selected.  The exact Qwen checkpoint revision, chat template, prompt text, GPU count, evaluation batch size, and non-placebo seeds are not given.

\findinglabel{Why it matters.} Dataset leakage and validation/test selection can inflate all reported effects, while missing selection and prompt details prevent independent reproduction.

\findinglabel{Suggested fix.} Add a dataset card, full category counts, expert-validation protocol, source/template-grouped split manifests, sampling seed, and licensing information.  Document the validation selection rule for layer, strength, and schedule; provide the exact model revision, prompts, execution configuration, and released split IDs.

\subsection*{8. Ablation, layer, and deployment-mechanism claims remain stronger than their evidence}

\classification{Unsupported interpretation and missing uncertainty}

\findinglabel{Location.} factorial ablation; Effect Sizes and Uncertainty; Layer-Wise Subspace Probing; Multi-Metric Trade-off \& Deployment Mitigation.

\findinglabel{Problem.} BERTScore differences between schedules are as small as 0.0001--0.0018, but no paired CIs, test statistics, or multiplicity policy are reported for the ablation.  All 12 rows again display exactly 7.32~s without dispersion.  The layer sweep reports only a range (0.7040--0.7140) and no per-layer uncertainty, yet infers that truthfulness is distributed across layers and that the framework is robust to layer shifts; a performance sweep cannot by itself establish either representation distribution or equivalence.

The Discussion further asserts a ``known tension'' in which steering forces attention onto clinical terminology templates and predicts that a repetition penalty or contrastive decoding can preserve the 5.20-point gain while restoring diversity.  Neither attention behavior nor the proposed mitigation is measured.  Positive steering changes RefPref by $+4.20$ points under bounded generation, whereas negative steering changes it by $-10.20$ points, so the stated symmetric-degradation purpose of the sign control is not realized or analyzed.

\findinglabel{Why it matters.} Small schedule ranks may be noise, and the representation, causal repetition, deployment, and symmetry conclusions do not follow from the reported summaries.

\findinglabel{Suggested fix.} Add paired ablation intervals and exploratory multiplicity handling, publish all per-layer results with uncertainty, and test layer equivalence rather than inferring it from a range.  Present the repetition mechanism and mitigation as hypotheses until directly evaluated, and analyze the positive/negative asymmetry.

\subsection*{9. Metric and hook details remain incomplete despite the improved environment list}

\classification{Implementation ambiguity}

\findinglabel{Location.} vector estimation; RefPref definition; Experimental Setup; Results tables.

\findinglabel{Problem.} The direction is extracted from the final token of a completed answer and injected into early autoregressive states without prompt-only, initial-token, pooled, or independent-error-template controls.  The exact hook module, zero- versus one-based layer convention, modified tensor positions, prompt-forward treatment, KV-cache behavior, and pre/post-hook logging order are still absent.  Rep-4 lacks a formula and tokenizer; Vietnamese ROUGE tokenization, BERTScore rescaling/baseline settings, and statistical software calls are not specified.  The placebo table's BERTScore 0.6945 is not identified as a mean, pooled score, or score from a particular direction.

\findinglabel{Why it matters.} These choices can materially change the vector, intervention, metrics, and statistical results even when package versions are fixed.

\findinglabel{Suggested fix.} Release executable hook pseudocode or code, exact tensor/index conventions, prompt and cache handling, complete metric configurations, tie counts, and statistical calls.  Report direction-level placebo metrics and test alternative extraction positions.

\section*{Equation, notation, sign, branch, and dimensional audit}

The direction $\mu_l^+-\mu_l^-$ and addition of $+\alpha(t)v_{\text{steer}}^{(l)}$ are sign-consistent with the intended positive intervention.  Unit normalization makes the vector dimensionally compatible with the hidden state when $\alpha(t)$ is an activation-scale coefficient.  The linear schedule and positive-$\alpha_0$ values of $D_{\text{continuous}}$, $D_{\text{cutoff}}$, and
\[
D_{\text{decay}}=\frac{K+1}{2}|\alpha_0|
\]
are arithmetically correct.  The reported positive matched coefficients 0.6375, 0.7650, and 0.8500 are also correct for $K=16$, $T=200$.  No inverse-trigonometric functions, coordinate transformations, or branch/quadrant choices occur.

The matched-dose definition remains sign-inconsistent as a general formula.  $D_{\text{decay}}$ is nonnegative, whereas
\[
\alpha_{\text{match}}=D_{\text{decay}}/T
=\frac{K+1}{2T}\alpha_0
\]
becomes negative for $\alpha_0<0$; the equality $\sum_t\alpha_{\text{match}}=D_{\text{decay}}$ then fails.  Restrict the control explicitly to $\alpha_0>0$, or define $\alpha_{\text{match}}=\operatorname{sgn}(\alpha_0)D_{\text{decay}}/T$ and match $\sum_t|\alpha_{\text{match}}|$.

The main notation drift is pre-hook $h_l^{(t)}$ versus post-hook $\hat h_l^{(t)}$ in the Abstract, contribution list, and Results.  The symbol $h_l(x_i,y_i)$ is first defined as the final-token representation of a completed prompt--answer sequence, whereas $h_l^{(t)}$ later denotes a generation-time residual stream; these roles should be distinguished explicitly.  $N$ is overloaded for training examples, category sizes, evaluation prompts, and random directions.  Token steps are discrete and should be written $t\in\{1,\ldots,100\}$ rather than $t\in[1,100]$.

\section*{Meaningful editorial, compilation, and reference findings}

\subsection*{The source compiles, but the natural-completion table severely overflows its column}

\findinglabel{Location.} Table~\texttt{tab:long\_gen}; clean two-pass pdfLaTeX log.

\findinglabel{Problem.} The prior environment mismatch and caption ampersand errors are fixed.  A clean two-pass build now succeeds, but the natural-completion table is approximately 61.44~pt wider than its one-column container.  The log also contains numerous high-badness underfull boxes, especially in the Abstract, contribution list, environment paragraph, and long technical paragraphs.

\findinglabel{Why it matters.} The natural table can intrude substantially into the adjacent IEEE column, and dense unbreakable content degrades typesetting quality.

\findinglabel{Suggested fix.} Use a legible \texttt{table*}, redesign the headings, or resize the table appropriately.  Break the software environment into compact prose or a reproducibility table, shorten long sentences, rebuild twice, and inspect the final PDF and log.

\subsection*{Citation fit remains uneven}

\findinglabel{Location.} Introduction; bibliography entries \texttt{ref10}, \texttt{ref11}, and \texttt{ref12}.

\findinglabel{Problem.} \texttt{ref10} is the LoRA paper rather than direct evidence that SFT mitigates the stated medical hallucinations.  \texttt{ref11} studies irrelevant retrieved context but does not establish that internal activation pathways cause evidence ignoring.  \texttt{ref12} concerns catastrophic forgetting generally rather than the asserted medical-language-model failure mode.

\findinglabel{Why it matters.} Direct evidence, general motivation, and mechanistic hypotheses are presented with similar certainty.

\findinglabel{Suggested fix.} Add directly matched sources or qualify these statements as motivation and hypotheses.

\section*{Proofing sweep}

The bundled mechanical scan was run on both the source and the clean compiled PDF.  Its only hit was the semicolon before ``Recall@3,'' which is not a capitalization defect.  Manual source, equation-adjacent, table, and bibliography checks identified these bounded corrections:

\begin{itemize}
    \item \textbf{Contribution~2 and natural Results:} replace the universal 100.00\% EOS statement with the condition-specific 99.80\%--100.00\% range.
    \item \textbf{Activation terminology:} use ``activation $\ell_2$ norm'' rather than ``Activation L2 norm,'' and use $\hat h_8^{(t)}$ consistently for post-hook measurements.
    \item \textbf{Timing units:} write values such as ``7.13~s'' rather than ``7.13s'' and define every $\pm$ term.
    \item \textbf{Natural table:} move the table to a full-width float or redesign it to eliminate the 61.44~pt overflow.
    \item \textbf{Abstract and Conclusion:} replace general ``accuracy,'' ``truthfulness,'' and ``confirms'' language with the measured endpoint and qualified mechanism wording.
\end{itemize}

The bibliography was checked for duplicated punctuation, malformed quotation punctuation, inconsistent capitalization, and obvious citation-format glitches; no additional mechanical defect was found.

\section*{Items requiring external verification}

\begin{itemize}
    \item Verify all formulary-derived references and synthetic counter-answers against the official source, current clinical guidance, qualified clinician review, and applicable reuse permissions.
    \item Validate multilingual BERTScore for Vietnamese clinical negation, quantities, units, contraindications, and interaction mechanisms before interpreting RefPref as factual alignment.
    \item Verify novelty against current token-dependent, scheduled, and dose-controlled activation-steering research.
    \item Verify the description of Qwen2.5-7B-Instruct as an SLM and the privacy/local-hospital claims using explicit definitions, a threat model, and measured resource requirements.
    \item Verify BM25 and stronger Vietnamese retrieval baselines before generalizing the RAG comparison.
\end{itemize}

\section*{Highest-priority revision sequence}

\begin{enumerate}
    \item Separate the bounded and natural paired analyses and publish every discordant count, test variant, CI procedure, and multiplicity calculation.
    \item Reframe the placebo result as finite-resolution direction-specific evidence rather than a calibrated five-sigma result.
    \item Align the Abstract and contribution claims with the limited activation-norm diagnostic, or add controlled multidimensional mechanism experiments.
    \item Add clinician-adjudicated correctness and harm endpoints, or narrow all clinical, truthfulness, and hallucination claims to RefPref.
    \item Present natural completion as a metric trade-off, resolve the Linear-Decay versus Hard-Cutoff method choice, and add paired multi-metric uncertainty.
    \item Strengthen and delimit the RAG comparison and provide component-wise timing and memory measurements.
    \item Complete dataset provenance, leakage controls, selection rules, metric/hook details, ablation uncertainty, and table layout.
\end{enumerate}

\end{document}